% !TEX root = paper.tex
\section{Related Work}\label{sec:related}

\subsection{%
	Diffusion Models for Image and Video Synthesis}\label{sub:related:diffusion_models}

Diffusion models
have emerged
as the dominant paradigm
for high-fidelity image generation
and increasingly constitute the foundation for video synthesis applications.
Foundational works establish generation
as iterative denoising\cite{%
	pmlr-v37-sohl-dickstein15, ho2020denoising},
providing strong baselines
for image quality benchmarks.
Score-based formulations
further connect diffusion
with denoising score matching
and establish continuous-time dynamics
\cite{%
	song2019generative},
while recent architectural innovations
push diffusion beyond GAN performance
on standard image benchmarks\cite{dhariwal2021diffusionmodelsbeatgans}.
For video synthesis,
diffusion models have been extended
to spatiotemporal settings,
achieving strong visual quality
through scalable attention mechanisms
and efficient broadcasting strategies
\cite{zhao2025realtimevideogenerationpyramid}.

\subsection{%
	Acceleration Methods for Diffusion Models}\label{sub:related:acceleration_methods}

Despite the remarkable generation quality
achieved by diffusion models
in image and video synthesis,
their inherently sequential inference procedure
and substantial computational requirements
pose significant barriers
to practical deployment
and real-time applications.
Existing approaches
to accelerating diffusion model inference
can be broadly categorized as follows:

\textbf{Sampling Acceleration.}
One fundamental approach to acceleration
reduces the number of denoising steps
via fast sampling algorithms,
including DDIM \cite{song2020denoising},
ODE/SDE solvers \cite{Karras2022edm,Lu_2025},
distillation-based methods \cite{salimans2022progressivedistillationfastsampling},
and consistency training \cite{song2023consistencymodels,luo2023latent}.

\textbf{Diffusion Quantization.}
Quantization compresses diffusion transformers
to reduce memory footprint and latency.
Post-training quantization (PTQ) methods
optimize deployment without retraining,
including PTQ4DiT \cite{wu2024ptq4dit},
which addresses outliers and timestep-wise distribution shifts,
enabling training-free low-bit quantization.
SVDQuant \cite{li2024svdquant}
uses low-rank branches to absorb outliers,
achieving 4-bit quantization with minimal quality loss.
For video models,
Q-VDiT \cite{feng2025q}
employs token-aware quantization and temporal distillation,
while ViDiT-Q \cite{zhao2024viditq}
adopts fine-grained dynamic quantization and mixed precision strategies.

\textbf{Feature Reuse.}
Given the computational overhead of training-dependent acceleration,
recent work has focused on training-free methods
that generalize across diverse diffusion models.
Feature reuse exploits activation redundancy across denoising steps
through intermediate caching and reuse,
offering superior generalizability.

Early approaches apply caching to U-Net-based models:
DeepCache \cite{Xu_2018} and Faster Diffusion \cite{li2024faster}
achieve acceleration through selective activation reuse.
Extending this to Diffusion Transformers,
FORA \cite{selvaraju2024forafastforwardcachingdiffusion},
ToCa \cite{zou2024accelerating},
DuCa \cite{zou2025rethinkingtokenwisefeaturecaching},
and $\Delta$-DiT \cite{chen2024deltadittrainingfreeaccelerationmethod}
employ heuristic-based policies
to cache attention and MLP layer activations.

Prediction-based methods take different approaches.
TeaCache \cite{liu2024timestep} utilizes carefully
 curated calibration data to fit polynomial coefficients,
while MagCache \cite{ma2025magcache}
achieves calibration with a single random sample.
TaylorSeer \cite{TaylorSeer2025} and HiCache \cite{feng2025hicachetrainingfreeaccelerationdiffusion}
employ polynomial-based extrapolation
to forecast feature evolution,
enabling prediction of future activations.

\textbf{Distinction from Prior Work:}
$D^3$-Cache fundamentally differs from existing feature reuse methods
through its foundation in Dynamic Mode Decomposition.
While prior approaches rely on heuristics or polynomial fitting,
$D^3$-Cache learns an optimal linear operator
via DMD to capture the underlying dynamics of transformer outputs.
This enables precise prediction of future transformer outputs
and adaptive skipping of redundant computations,
achieving superior quality-efficiency trade-offs
compared to existing cache-based methods.
